% macros.tex
% Late-binding citation macros for the dissertation paper.
%
% Every number in prose resolves at build time through these macros.
% Values are defined in generated/values.tex via \csname lookup;
% build_numbers.py produces that file from data/*.csv and
% generated/figures/*.numbers.csv.
%
% Namespace prefix: ds@ (dissertation). Must match NAMESPACE in build_numbers.py.
%
% Usage:
%   \result{csv_stem}{key}[precision]
%   \resdelta{csv_stem}{key_a}{key_b}[precision]
%   \resratio{csv_stem}{key_a}{key_b}[precision]
%   \resultCI{csv_stem}{mean_key,lo_key,hi_key}[precision]
%   \resultPM{csv_stem}{mean_key,err_key}[precision]
%
% Where:
%   csv_stem  - stem of the CSV file (e.g. "headline" for headline.csv,
%               "fig.02_signflip" for the sidecar of Figure 2 — note the
%               filename is still 02_signflip.* for build-pipeline stability
%               even though the §5.3 prose framing was retuned 2026-05-28;
%               see paper/05_results.md §5.3 for the rationale)
%   key       - colon-joined coordinate per the CSV's key_schema
%   precision - number of decimal places for siunitx rounding (optional;
%               default is the CSV's default_precision, embedded as
%               ds@<stem>@__default_precision__ in values.tex)

\usepackage{siunitx}
\usepackage{xfp}
\usepackage{etoolbox}

% --------------------------------------------------------------------------
% Definitions below use @ in macro names. Outside a .sty file, @ has catcode
% 12, so we temporarily make it a letter via \makeatletter / \makeatother.
% --------------------------------------------------------------------------
\makeatletter

% --------------------------------------------------------------------------
% Internal helper: look up \csname ds@<stem>@<key>\endcsname
% Raises \PackageError if undefined.
% --------------------------------------------------------------------------
\newcommand{\ocget}[2]{%
  \@ifundefined{ds@#1@#2}{%
    \PackageError{dissertation-macros}%
      {Undefined result key '#2' in CSV stem '#1'.
       Check data/#1.csv and re-run 'make values'.}%
      {The key '#2' was not found in the values generated from '#1'.
       Valid keys are listed in generated/values.tex.}%
    \textbf{??}%
  }{%
    \csname ds@#1@#2\endcsname%
  }%
}

% --------------------------------------------------------------------------
% Helper: get default precision for a CSV stem (falls back to 3)
% --------------------------------------------------------------------------
\newcommand{\ocgetdefprec}[1]{%
  \@ifundefined{ds@#1@__default_precision__}{3}{\csname ds@#1@__default_precision__\endcsname}%
}

% --------------------------------------------------------------------------
% \result{csv}{key}[precision]
% --------------------------------------------------------------------------
\DeclareRobustCommand{\result}[2]{%
  \@ifnextchar[%]
    {\ds@result@prec{#1}{#2}}%
    {\ds@result@prec{#1}{#2}[\ocgetdefprec{#1}]}%
}
\def\ds@result@prec#1#2[#3]{%
  \num[round-mode=places,round-precision=#3]{\ocget{#1}{#2}}%
}

% --------------------------------------------------------------------------
% \resdollar{csv}{key}[precision]
% Render a USD value with the sign OUTSIDE the dollar sign: a negative
% value prints as "-\$0.023", not "\$-0.023". Rounds first so that values
% that round to zero print "\$0.000" rather than "-\$0.000". Use inside
% math mode (the bare "-" is a math minus), mirroring the old
% "$\$\result{...}$" idiom that placed the minus after the dollar sign.
% --------------------------------------------------------------------------
\DeclareRobustCommand{\resdollar}[2]{%
  \@ifnextchar[%]
    {\ds@dollar@prec{#1}{#2}}%
    {\ds@dollar@prec{#1}{#2}[\ocgetdefprec{#1}]}%
}
\def\ds@dollar@prec#1#2[#3]{%
  \edef\ds@dv{\fpeval{round(\ocget{#1}{#2},#3)}}%
  \ifnum\fpeval{\ds@dv<0}=1 %
    -\$\num[round-mode=places,round-precision=#3]{\fpeval{abs(\ds@dv)}}%
  \else
    \$\num[round-mode=places,round-precision=#3]{\ds@dv}%
  \fi
}

% --------------------------------------------------------------------------
% \resdelta{csv}{key_a}{key_b}[precision]
% Computes key_a - key_b using \fpeval from xfp.
% (Named \resdelta, not \delta, to avoid clashing with the Greek letter.)
% --------------------------------------------------------------------------
\DeclareRobustCommand{\resdelta}[3]{%
  \@ifnextchar[%]
    {\ds@delta@prec{#1}{#2}{#3}}%
    {\ds@delta@prec{#1}{#2}{#3}[\ocgetdefprec{#1}]}%
}
\def\ds@delta@prec#1#2#3[#4]{%
  \edef\ds@va{\ocget{#1}{#2}}%
  \edef\ds@vb{\ocget{#1}{#3}}%
  \num[round-mode=places,round-precision=#4,explicit-sign]{\fpeval{\ds@va - \ds@vb}}%
}

% --------------------------------------------------------------------------
% \resratio{csv}{key_a}{key_b}[precision]
% Computes key_a / key_b using \fpeval.
% --------------------------------------------------------------------------
\DeclareRobustCommand{\resratio}[3]{%
  \@ifnextchar[%]
    {\ds@ratio@prec{#1}{#2}{#3}}%
    {\ds@ratio@prec{#1}{#2}{#3}[\ocgetdefprec{#1}]}%
}
\def\ds@ratio@prec#1#2#3[#4]{%
  \edef\ds@va{\ocget{#1}{#2}}%
  \edef\ds@vb{\ocget{#1}{#3}}%
  \num[round-mode=places,round-precision=#4]{\fpeval{\ds@va / \ds@vb}}%
}

% --------------------------------------------------------------------------
% \resp{csv}{key}[sigfigs]
% Render a p-value at N significant figures (default 2). Uses siunitx
% scientific notation automatically for small p (e.g. 2.7e-17 → 2.7×10⁻¹⁷),
% which `\result`'s round-mode=places would clobber to 0.000.
% --------------------------------------------------------------------------
\DeclareRobustCommand{\resp}[2]{%
  \@ifnextchar[%]
    {\ds@resp@prec{#1}{#2}}%
    {\ds@resp@prec{#1}{#2}[2]}%
}
\def\ds@resp@prec#1#2[#3]{%
  \num[round-mode=figures,round-precision=#3]{\ocget{#1}{#2}}%
}

% --------------------------------------------------------------------------
% \respp{csv}{key}[precision]
% Render a stored 0–1 fraction as percentage points: value × 100. Pass-rate
% deltas in paired_contrasts.csv are fractions; tables typically want pp.
% Default precision: 2 decimal places.
% --------------------------------------------------------------------------
\DeclareRobustCommand{\respp}[2]{%
  \@ifnextchar[%]
    {\ds@respp@prec{#1}{#2}}%
    {\ds@respp@prec{#1}{#2}[2]}%
}
\def\ds@respp@prec#1#2[#3]{%
  \num[round-mode=places,round-precision=#3]{\fpeval{100*\ocget{#1}{#2}}}%
}

% --------------------------------------------------------------------------
% \respct{csv}{delta_key}{ref_key}[precision]
% Render a paired-contrast delta as a percentage of a reference mean:
% (delta / ref) × 100, with a literal '%' appended. For "code-only is X%
% cheaper than competitor" framing, pass the contrast's mean_b cell as ref.
% Default precision: 1 decimal place.
% --------------------------------------------------------------------------
\DeclareRobustCommand{\respct}[3]{%
  \@ifnextchar[%]
    {\ds@respct@prec{#1}{#2}{#3}}%
    {\ds@respct@prec{#1}{#2}{#3}[1]}%
}
\def\ds@respct@prec#1#2#3[#4]{%
  \edef\ds@vd{\ocget{#1}{#2}}%
  \edef\ds@vr{\ocget{#1}{#3}}%
  \num[round-mode=places,round-precision=#4]{\fpeval{100*\ds@vd/\ds@vr}}\%%
}

% --------------------------------------------------------------------------
% \resultCI{csv}{mean_key,lo_key,hi_key}[precision]
% Renders as: mean [lo, hi]
% --------------------------------------------------------------------------
\DeclareRobustCommand{\resultCI}[2]{%
  \@ifnextchar[%]
    {\ds@ci@prec{#1}{#2}}%
    {\ds@ci@prec{#1}{#2}[\ocgetdefprec{#1}]}%
}
\def\ds@ci@prec#1#2[#3]{%
  \ds@splitthree{#2}{\ds@keyA}{\ds@keyB}{\ds@keyC}%
  \num[round-mode=places,round-precision=#3]{\ocget{#1}{\ds@keyA}}%
  \ [\num[round-mode=places,round-precision=#3]{\ocget{#1}{\ds@keyB}},
     \num[round-mode=places,round-precision=#3]{\ocget{#1}{\ds@keyC}}]%
}

% --------------------------------------------------------------------------
% \resultPM{csv}{mean_key,err_key}[precision]
% Renders as: mean ± err
% --------------------------------------------------------------------------
\DeclareRobustCommand{\resultPM}[2]{%
  \@ifnextchar[%]
    {\ds@pm@prec{#1}{#2}}%
    {\ds@pm@prec{#1}{#2}[\ocgetdefprec{#1}]}%
}
\def\ds@pm@prec#1#2[#3]{%
  \ds@splittwo{#2}{\ds@keyA}{\ds@keyB}%
  \num[round-mode=places,round-precision=#3]{\ocget{#1}{\ds@keyA}}%
  $\pm$%
  \num[round-mode=places,round-precision=#3]{\ocget{#1}{\ds@keyB}}%
}

% --------------------------------------------------------------------------
% Utility: split comma-delimited pairs and triples into named macros.
% --------------------------------------------------------------------------
\def\ds@splittwo#1#2#3{%
  \ds@splittwoaux#1\@nil{#2}{#3}%
}
\def\ds@splittwoaux#1,#2\@nil#3#4{%
  \def#3{#1}\def#4{#2}%
}

\def\ds@splitthree#1#2#3#4{%
  \ds@splitthreeaux#1\@nil{#2}{#3}{#4}%
}
\def\ds@splitthreeaux#1,#2,#3\@nil#4#5#6{%
  \def#4{#1}\def#5{#2}\def#6{#3}%
}

\makeatother
